\documentclass[12pt, a4paper]{article}

% ==========================================
%  Packages Setup
% ==========================================
\usepackage[utf8]{inputenc}
\usepackage{amsmath, amssymb, amsthm}
\usepackage{geometry}
\geometry{top=2.5cm, bottom=2.5cm, left=2.5cm, right=2.5cm}
\usepackage{setspace}
\onehalfspacing
\usepackage{hyperref}
\hypersetup{
    colorlinks=true,
    linkcolor=black,
    urlcolor=blue
}
\usepackage{graphicx}
\usepackage{float}
\usepackage{xcolor}

% ==========================================
%  Code Listing Setup
% ==========================================
\usepackage{listings}
\lstset{
    language=Python,
    backgroundcolor=\color{black!5},
    basicstyle=\ttfamily\small,
    keywordstyle=\color{blue}\bfseries,
    stringstyle=\color{red},
    commentstyle=\color{green!60!black},
    numbers=left,
    numberstyle=\tiny\color{gray},
    stepnumber=1,
    frame=single,
    rulecolor=\color{black!20},
    breaklines=true
}

% ==========================================
%  Theorem Environments
% ==========================================
\newtheoremstyle{break}
  {\topsep}
  {\topsep}
  {\normalfont}
  {}
  {\bfseries}
  {.}
  {\newline}
  {}

\theoremstyle{break}
\newtheorem{theorem}{Theorem}[section]
\newtheorem{definition}{Definition}[section]
\newtheorem{corollary}{Corollary}[theorem]
\renewcommand{\qedsymbol}{$\blacksquare$}

\makeatletter
\renewenvironment{proof}[1][\proofname]{\par
  \pushQED{\qed}%
  \normalfont \topsep6\p@\@plus6\p@\relax
  \trivlist
  \item[\hskip\labelsep\bfseries #1\@addpunct{.}]\mbox{}\newline\ignorespaces
}{%
  \popQED\endtrivlist\@endpefalse
}
\makeatother
\providecommand{\tightlist}{\setlength{\itemsep}{0pt}\setlength{\parskip}{0pt}}
\providecommand{\passthrough}[1]{#1}

\begin{document}

% ==========================================
%  Title Page
% ==========================================
\begin{titlepage}
    \begin{center}
        \vspace*{2cm}
        {\LARGE \bfseries Lecture Notes on Scientific Computing \par}
        \vspace{3cm}
        \begin{tabular}{rl}
            \textbf{Author:} & Windson \\
            \textbf{Course:} & Scientific Computing \\
        \end{tabular}
        \vspace{1.5cm}

        {\large Department of Applied Mathematics \par}
        {\large National Chung Hsing University \par}
        \vspace{0.3cm}
        145 Xingda Road, South District, Taichung City 40227, Taiwan (R.O.C.) \par
        e-mail: info AT windson.cc

        \vfill
        {\large \textbf{License} \par}
        {\small This document is licensed under the Creative Commons Attribution 4.0 International License (CC BY-SA 4.0).\par}
        \vspace{1cm}
        {\large March 2, 2026 \par}
    \end{center}
\end{titlepage}

% ==========================================
%  Table of Contents
% ==========================================
\pagenumbering{roman}
\tableofcontents
\clearpage
\pagenumbering{arabic}
\setcounter{page}{1}

% ==========================================
%  Main Content
Scientific computing develops numerical methods for problems whose exact solutions are unavailable, inconvenient, or too expensive to evaluate. The central workflow is:

\begin{enumerate}
\def\labelenumi{\arabic{enumi}.}
\tightlist
\item
  formulate a mathematical model;
\item
  replace derivatives, integrals, or functions by computable approximations;
\item
  derive an error estimate;
\item
  study convergence and stability; and
\item
  validate the method with computation.
\end{enumerate}

This note collects the interpolation, numerical integration, initial-value problem, Runge--Kutta, stability, and nonlinear-pendulum material from the course.

\section{Polynomial Interpolation}\label{polynomial-interpolation}

Suppose that a function is known at distinct nodes

\[
(x_0,y_0),(x_1,y_1),\ldots,(x_n,y_n), \qquad y_i=f(x_i).
\]

We seek a polynomial \(P_n\) of degree at most \(n\) satisfying

\[
P_n(x_i)=y_i, \qquad i=0,1,\ldots,n.
\]

\subsection{Lagrange Basis}\label{lagrange-basis}

For each node \(x_i\), construct a polynomial that vanishes at every other node. Its unnormalized numerator is

\[
\prod_{\substack{j=0\\j\ne i}}^n (x-x_j).
\]

At \(x=x_i\), this numerator has the value

\[
\prod_{\substack{j=0\\j\ne i}}^n (x_i-x_j).
\]

Dividing by this value normalizes the polynomial to one:

\[
L_i(x)=\prod_{\substack{j=0\\j\ne i}}^n
\frac{x-x_j}{x_i-x_j}.
\]

Indeed, if \(k\ne i\), one factor is \(x_k-x_k=0\), so \(L_i(x_k)=0\). If \(k=i\), every numerator and denominator factor is equal, so \(L_i(x_i)=1\). Therefore,

\[
L_i(x_j)=\delta_{ij}=\begin{cases}1,&i=j,\\0,&i\ne j.\end{cases}
\]

Taking a linear combination with weights \(y_i\) gives the Lagrange interpolating polynomial:

\[
P_n(x)=\sum_{i=0}^n y_iL_i(x)
=\sum_{i=0}^n y_i\prod_{\substack{j=0\\j\ne i}}^n
\frac{x-x_j}{x_i-x_j}.
\]

At a node \(x_k\),

\[
P_n(x_k)=\sum_{i=0}^n y_iL_i(x_k)
=\sum_{i=0}^n y_i\delta_{ik}=y_k,
\]

so the polynomial interpolates every data point.

\subsection{Example}\label{example}

For \((x_0,y_0)=(1,2)\), \((x_1,y_1)=(2,4)\), and \((x_2,y_2)=(4,8)\), the basis function associated with \(x_1=2\) is

\[
\begin{aligned}
L_1(x)&=\frac{(x-x_0)(x-x_2)}{(x_1-x_0)(x_1-x_2)}\\
&=\frac{(x-1)(x-4)}{(2-1)(2-4)}\\
&=-\frac12(x-1)(x-4).
\end{aligned}
\]

At \(x=3\),

\[
L_0(3)=-\frac13,\qquad L_1(3)=1,\qquad L_2(3)=\frac13.
\]

Consequently,

\[
P_2(3)=2\left(-\frac13\right)+4(1)+8\left(\frac13\right)=6.
\]

\subsection{Uniqueness}\label{uniqueness}

Suppose \(P_n\) and \(Q_n\) are two polynomials of degree at most \(n\) that interpolate the same \(n+1\) data points. Let

\[
R(x)=P_n(x)-Q_n(x).
\]

Then \(R(x_i)=0\) for every \(i\), so \(R\) has at least \(n+1\) distinct roots. A nonzero polynomial of degree at most \(n\) can have at most \(n\) roots. Hence \(R\) is the zero polynomial and

\[
P_n(x)=Q_n(x).
\]

The interpolating polynomial is therefore unique.

\subsection{Partition of Unity and Affine Invariance}\label{partition-of-unity-and-affine-invariance}

Apply interpolation to the constant function \(f(x)=1\). Since \(f(x_i)=1\),

\[
1=P_n(x)=\sum_{i=0}^n L_i(x).
\]

Thus the basis functions form a partition of unity: their weights always sum to one. In particular, interpolation preserves constants and affine combinations.

If the independent variable is changed by an affine map \(x=at+b\), the same construction in the \(t\)-coordinate produces the correspondingly transformed interpolating polynomial. Translation and scaling change the coordinates, but not the interpolation structure.

\subsection{Nodal Polynomial Form}\label{nodal-polynomial-form}

Define the nodal polynomial

\[
Q(x)=\prod_{j=0}^n(x-x_j).
\]

Removing the factor corresponding to \(x_i\) gives

\[
\frac{Q(x)}{x-x_i}=\prod_{\substack{j=0\\j\ne i}}^n(x-x_j).
\]

Differentiating the product and evaluating at \(x_i\), every term vanishes except the term in which the factor \((x-x_i)\) is differentiated:

\[
Q'(x_i)=\prod_{\substack{j=0\\j\ne i}}^n(x_i-x_j).
\]

Therefore,

\[
L_i(x)=\frac{Q(x)}{(x-x_i)Q'(x_i)},
\]

and the interpolant can also be written as

\[
P_n(x)=\sum_{i=0}^n
\frac{Q(x)}{(x-x_i)Q'(x_i)}f(x_i).
\]

\subsection{Interpolation Error}\label{interpolation-error}

Assume \(f\in C^{n+1}[a,b]\), with all nodes and the evaluation point \(x\) in the interval. Define

\[
E(x)=f(x)-P_n(x).
\]

At the interpolation nodes, \(E(x_i)=0\). For a point \(x\) not equal to a node, choose a coefficient \(c(x)\) so that

\[
c(x)=\frac{f(x)-P_n(x)}{Q(x)}.
\]

Now define an auxiliary function of \(z\):

\[
F(z)=f(z)-P_n(z)-c(x)Q(z).
\]

By construction, \(F(x_i)=0\) for \(i=0,\ldots,n\), and \(F(x)=0\). Thus \(F\) has \(n+2\) distinct roots. Repeated use of Rolle's theorem gives a point \(\xi\) between the nodes and \(x\) such that

\[
F^{(n+1)}(\xi)=0.
\]

Since \(P_n^{(n+1)}=0\) and \(Q^{(n+1)}=(n+1)!\),

\[
0=f^{(n+1)}(\xi)-c(x)(n+1)!.
\]

Therefore,

\[
c(x)=\frac{f^{(n+1)}(\xi)}{(n+1)!},
\]

and the error formula is

\[
\boxed{
f(x)-P_n(x)=\frac{f^{(n+1)}(\xi)}{(n+1)!}
\prod_{i=0}^n(x-x_i)
}.
\]

If \(f\) is a polynomial of degree \(k\le n\), then \(f^{(n+1)}=0\). Hence the interpolation is exact: \(E(x)=0\) for every \(x\).

\subsection{Derivatives of the Interpolant}\label{derivatives-of-the-interpolant}

Differentiate the interpolant term by term:

\[
P_n'(x)=\sum_{i=0}^n f(x_i)L_i'(x).
\]

Starting with

\[
Q(x)=(x-x_i)Q'(x_i)L_i(x),
\]

differentiate once:

\[
Q'(x)=Q'(x_i)L_i(x)+(x-x_i)Q'(x_i)L_i'(x).
\]

At another node \(x_j\), \(j\ne i\), \(L_i(x_j)=0\), so

\[
L_i'(x_j)=\frac{Q'(x_j)}{(x_j-x_i)Q'(x_i)}.
\]

At \(x=x_i\), differentiate the identity once more:

\[
Q''(x)=2Q'(x_i)L_i'(x)+(x-x_i)Q'(x_i)L_i''(x).
\]

Evaluating at \(x_i\) gives

\[
L_i'(x_i)=\frac{Q''(x_i)}{2Q'(x_i)}.
\]

The derivative error follows by differentiating the interpolation error. Writing the intermediate point as \(\xi=\xi(x)\),

\[
f'(x)-P_n'(x)=\frac{1}{(n+1)!}
\left[Q'(x)f^{(n+1)}(\xi(x))
+Q(x)f^{(n+2)}(\xi(x))\xi'(x)\right].
\]

At a node \(x_i\), the second term disappears because \(Q(x_i)=0\). Hence

\[
f'(x_i)-P_n'(x_i)=
\frac{f^{(n+1)}(\xi_i)}{(n+1)!}
\prod_{\substack{k=0\\k\ne i}}^n(x_i-x_k).
\]

\subsection{Runge's Phenomenon}\label{runges-phenomenon}

Increasing the number of nodes does not always improve interpolation. For

\[
f(x)=\frac{1}{1+25x^2}, \qquad -1\le x\le 1,
\]

use equally spaced nodes \(x_i=-1+2i/n\). The unique polynomial is

\[
P_n(x)=\sum_{i=0}^n f(x_i)
\prod_{\substack{j=0\\j\ne i}}^n\frac{x-x_j}{x_i-x_j}.
\]

For large \(n\), the product factors become badly conditioned near the endpoints and the interpolant can oscillate strongly there. This is Runge's phenomenon. More nodes are not sufficient by themselves; node placement, such as Chebyshev nodes, is also important.

\section{Numerical Integration}\label{numerical-integration}

Let

\[
I=\int_a^b f(x)\,dx,
\]

and divide \([a,b]\) into \(n\) equal subintervals. Set

\[
h=\Delta x=\frac{b-a}{n},\qquad x_i=a+ih.
\]

\subsection{Left and Right Rectangle Rules}\label{left-and-right-rectangle-rules}

The left-point approximation is

\[
I_L=h\sum_{i=1}^n f(x_{i-1}),
\]

and the right-point approximation is

\[
I_R=h\sum_{i=1}^n f(x_i).
\]

For the right rule, write the error as a sum of local errors:

\[
\begin{aligned}
E_R&=I-I_R\\
&=\sum_{i=1}^n\int_{x_{i-1}}^{x_i}[f(x)-f(x_i)]\,dx.
\end{aligned}
\]

Taylor's theorem about the right endpoint gives

\[
f(x)=f(x_i)+f'(x_i)(x-x_i)+\frac{f''(\xi_i)}{2}(x-x_i)^2.
\]

Therefore,

\[
E_R=\sum_{i=1}^n\int_{x_{i-1}}^{x_i}
\left[f'(x_i)(x-x_i)+\frac{f''(\xi_i)}2(x-x_i)^2\right]dx.
\]

If \(M_1=\max_{[a,b]}|f'(x)|\), the leading term satisfies

\[
\left|\int_{x_{i-1}}^{x_i}f'(x_i)(x-x_i)\,dx\right|
\le \frac{M_1h^2}{2}.
\]

Summing \(n\) intervals,

\[
|E_R|\lesssim n\frac{M_1h^2}{2}
=\frac{M_1(b-a)^2}{2n}.
\]

The same argument about the left endpoint gives

\[
|E_L|\lesssim\frac{M_1(b-a)^2}{2n}.
\]

Thus both endpoint rules are first-order methods: their global error is \(O(h)=O(n^{-1})\).

\subsection{Midpoint Rule}\label{midpoint-rule}

Let \(m_i=(x_{i-1}+x_i)/2\). Approximate \(f\) about the midpoint:

\[
f(x)=f(m_i)+f'(m_i)(x-m_i)+\frac{f''(\xi_i)}2(x-m_i)^2.
\]

The local midpoint error is

\[
E_{M,i}=\int_{x_{i-1}}^{x_i}[f(x)-f(m_i)]\,dx.
\]

The linear term cancels by symmetry:

\[
\int_{x_{i-1}}^{x_i}(x-m_i)\,dx=0.
\]

For the quadratic term, use \(t=x-m_i\), with limits \(-h/2\) and \(h/2\):

\[
\int_{-h/2}^{h/2}\frac{f''(\xi_i)}2t^2\,dt
=\frac{f''(\xi_i)}2\frac{h^3}{12}
=\frac{f''(\xi_i)h^3}{24}.
\]

If \(M_2=\max|f''|\), then

\[
|E_M|\le n\frac{M_2h^3}{24}
=\frac{M_2(b-a)^3}{24n^2}.
\]

The midpoint rule is second-order.

\subsection{Trapezoidal Rule}\label{trapezoidal-rule}

On \([x_{i-1},x_i]\), the trapezoidal approximation is

\[
T_i=\frac h2[f(x_{i-1})+f(x_i)].
\]

Summing gives

\[
T_n=h\left[\frac{f(x_0)}2+\sum_{i=1}^{n-1}f(x_i)+\frac{f(x_n)}2\right].
\]

The local trapezoidal error is

\[
E_{T,i}=-\frac{h^3}{12}f''(\xi_i),
\]

so

\[
|E_T|\le\frac{M_2(b-a)^3}{12n^2}.
\]

The trapezoidal rule is also second-order.

\subsection{Simpson's Rule from Lagrange Interpolation}\label{simpsons-rule-from-lagrange-interpolation}

Use three consecutive points \(x_{i-1}\), \(x_i\), and \(x_{i+1}\), where the spacing is \(h\). The quadratic Lagrange interpolant is

\[
P_2(x)=L_0(x)f(x_{i-1})+L_1(x)f(x_i)+L_2(x)f(x_{i+1}),
\]

where

\[
\begin{aligned}
L_0(x)&=\frac{(x-x_i)(x-x_{i+1})}{2h^2},\\
L_1(x)&=-\frac{(x-x_{i-1})(x-x_{i+1})}{h^2},\\
L_2(x)&=\frac{(x-x_{i-1})(x-x_i)}{2h^2}.
\end{aligned}
\]

Set \(x=x_i+ht\), so \(dx=h\,dt\) and \(t\in[-1,1]\). Then

\[
\int_{x_{i-1}}^{x_{i+1}}L_0(x)\,dx=\frac h3,
\]

because

\[
\frac h2\int_{-1}^1(t^2-t)\,dt=\frac h3.
\]

Similarly,

\[
\int_{x_{i-1}}^{x_{i+1}}L_1(x)\,dx=\frac{4h}{3},
\]

because

\[
-h\int_{-1}^1(t^2-1)\,dt=\frac{4h}{3},
\]

and

\[
\int_{x_{i-1}}^{x_{i+1}}L_2(x)\,dx=\frac h3.
\]

Therefore,

\[
\int_{x_{i-1}}^{x_{i+1}}f(x)\,dx
\approx\frac h3[f(x_{i-1})+4f(x_i)+f(x_{i+1})].
\]

Adding the \(n/2\) panels, with \(n\) even, gives the composite Simpson rule:

\[
S_n=\frac h3\left[f(x_0)+4f(x_1)+2f(x_2)+4f(x_3)+\cdots+4f(x_{n-1})+f(x_n)\right].
\]

Taylor expansion about the midpoint shows that the odd powers cancel. The first remaining neglected term is proportional to \(h^5f^{(4)}\), so the global error is \(O(h^4)\).

\subsection{Numerical Experiment}\label{numerical-experiment}

For a smooth test function such as \(f(x)=e^{-x^2}\), compute a high-precision reference value and compare the five rules for \(N=10,20,40,80,160\). If \(E_N\) denotes the error, the observed order is estimated by

\[
p\approx\frac{\log(|E_N|/|E_{2N}|)}{\log 2}.
\]

The expected orders are one for left/right rules, two for midpoint and trapezoidal rules, and four for Simpson's rule. Symmetry can occasionally produce a higher observed order for a particular function, but the general theoretical guarantees remain those listed above.

\section{Initial-Value Problems}\label{initial-value-problems}

Consider

\[
u'(t)=f(t,u(t)),\qquad u(0)=u_0.
\]

Integrating from \(t\) to \(t+h\) and applying the fundamental theorem of calculus gives the exact identity

\[
u(t+h)=u(t)+\int_t^{t+h}f(s,u(s))\,ds.
\]

Every one-step method below approximates this integral in a different way.

\subsection{Euler Method}\label{euler-method}

Approximate the integrand by its left endpoint:

\[
f(s,u(s))\approx f(t,u(t)).
\]

Then

\[
\int_t^{t+h}f(s,u(s))\,ds\approx hf(t,u(t)),
\]

and, with \(t_n=nh\) and \(v_n\approx u(t_n)\),

\[
\boxed{v_{n+1}=v_n+hf(t_n,v_n)}.
\]

Writing \(K_1=f(t_n,v_n)\) gives \(v_{n+1}=v_n+hK_1\).

\subsection{Midpoint Method}\label{midpoint-method}

Approximate the integral with the slope at the midpoint:

\[
u(t+h)\approx u(t)+hf\left(t+\frac h2,u\left(t+\frac h2\right)\right).
\]

The midpoint value is unknown, so predict it by a half Euler step:

\[
v^{(1)}=v_n+\frac h2f(t_n,v_n).
\]

Thus

\[
\begin{aligned}
K_1&=f(t_n,v_n),\\
K_2&=f\left(t_n+\frac h2,v_n+\frac h2K_1\right),\\
v_{n+1}&=v_n+hK_2.
\end{aligned}
\]

\subsection{Heun's Method}\label{heuns-method}

Apply the trapezoidal rule to the IVP integral:

\[
u(t+h)\approx u(t)+\frac h2[f(t,u(t))+f(t+h,u(t+h))].
\]

Predict the unknown endpoint by Euler:

\[
v^{(1)}=v_n+hf(t_n,v_n).
\]

Then

\[
\begin{aligned}
K_1&=f(t_n,v_n),\\
K_2&=f(t_n+h,v_n+hK_1),\\
v_{n+1}&=v_n+\frac h2(K_1+K_2).
\end{aligned}
\]

Both midpoint and Heun methods have local error \(O(h^3)\) and global error \(O(h^2)\).

\subsection{Kutta's Third-Order Method}\label{kuttas-third-order-method}

One third-order method from the course is

\[
\begin{aligned}
K_1&=f(t_n,v_n),\\
K_2&=f(t_n+h,v_n+hK_1),\\
K_3&=f\left(t_n+\frac h2,v_n+\frac h2\frac{K_1+K_2}{2}\right),\\
v_{n+1}&=v_n+h\frac{K_1+K_2+4K_3}{6}.
\end{aligned}
\]

The weights are chosen so that the Taylor expansion agrees through terms of order three.

\subsection{Classical RK4}\label{classical-rk4}

Use two midpoint estimates, one based on \(K_1\) and one based on \(K_2\):

\[
\begin{aligned}
K_1&=f(t_n,v_n),\\
K_2&=f\left(t_n+\frac h2,v_n+\frac h2K_1\right),\\
K_3&=f\left(t_n+\frac h2,v_n+\frac h2K_2\right),\\
K_4&=f(t_n+h,v_n+hK_3).
\end{aligned}
\]

The update is

\[
\boxed{v_{n+1}=v_n+\frac h6(K_1+2K_2+2K_3+K_4)}.
\]

The four slopes approximate the integral with a Simpson-like weighted average. Taylor matching gives local truncation error \(O(h^5)\) and global error \(O(h^4)\).

\section{Error, Consistency, and Stability}\label{error-consistency-and-stability}

\subsection{ Euler Residual}\label{euler-residual}

Let \(u_n=u(t_n)\) be the exact solution and \(v_n\) the numerical solution. The residual of Euler's scheme is

\[
r_n=\frac{u_{n+1}-u_n}{h}-f(t_n,u_n).
\]

Taylor's theorem gives

\[
u_{n+1}=u_n+hu'(t_n)+\frac{h^2}{2}u''(\xi_n),
\]

so, because \(u'=f(t,u)\),

\[
r_n=\frac h2u''(\xi_n)=O(h).
\]

Since \(r_n\to0\) as \(h\to0\), the method is consistent and first-order accurate as an approximation to the differential equation.

\subsection{Euler Error Equation}\label{euler-error-equation}

Let \(e_n=u_n-v_n\). Subtract the numerical update from the exact equation with residual:

\[
\frac{u_{n+1}-u_n}{h}=f(t_n,u_n)+r_n,
\qquad
\frac{v_{n+1}-v_n}{h}=f(t_n,v_n).
\]

Then

\[
e_{n+1}=e_n+h[f(t_n,u_n)-f(t_n,v_n)]+hr_n.
\]

By the mean-value theorem,

\[
f(t_n,u_n)-f(t_n,v_n)=f_u(t_n,\eta_n)e_n,
\]

where \(\eta_n\) lies between \(u_n\) and \(v_n\). Hence

\[
e_{n+1}=[1+hf_u(t_n,\eta_n)]e_n+hr_n.
\]

Assume \(|u''(t)|\le M\) and \(|f_u(t,u)|\le L\). Then

\[
|e_{n+1}|\le(1+hL)|e_n|+\frac{h^2M}{2}.
\]

Iterating from \(e_0=0\),

\[
|e_{n+1}|\le\frac{h^2M}{2}\sum_{k=0}^n(1+hL)^k
=\frac{hM}{2L}[(1+hL)^{n+1}-1].
\]

Since \(1+hL\le e^{hL}\) and \(T=(n+1)h\),

\[
|e_{n+1}|\le\frac{hM}{2L}(e^{LT}-1).
\]

For fixed \(T\), this tends to zero linearly with \(h\).

\subsection{Stability and Amplification Factors}\label{stability-and-amplification-factors}

For the test equation

\[
u'=\lambda u,
\]

Euler gives

\[
v_{n+1}=(1+h\lambda)v_n.
\]

The amplification factor is

\[
R(z)=1+z,\qquad z=h\lambda.
\]

Absolute stability requires

\[
|R(z)|=|1+h\lambda|\le1.
\]

For \(\lambda<0\), this becomes

\[
-1\le1+h\lambda\le1,
\]

so

\[
0\le h\le\frac{2}{|\lambda|}.
\]

For example, \(u'=-5u\) requires \(h\le0.4\) for stable explicit Euler integration. For a general Runge--Kutta method, the stability condition is \(|R(h\lambda)|\le1\); for a system, \(h\lambda_i\) must lie in the method's stability region for every relevant Jacobian eigenvalue \(\lambda_i\).

\section{Nonlinear Simple Pendulum}\label{nonlinear-simple-pendulum}

\subsection{Physical Model}\label{physical-model}

Consider a point mass \(m\) attached to a massless, inextensible string of length \(L\) in a uniform gravitational field \(g\). Let \(\theta(t)\) be the angle measured from the downward vertical. The arc length is

\[
s=L\theta.
\]

The gravitational force has tangential component \(-mg\sin\theta\). Newton's second law in the tangential direction gives

\[
-mg\sin\theta=m\frac{d^2s}{dt^2}.
\]

Since \(s=L\theta\),

\[
\frac{d^2s}{dt^2}=L\frac{d^2\theta}{dt^2}.
\]

Substituting and cancelling \(m\) yields

\[
\boxed{\theta''+\frac gL\sin\theta=0}.
\]

The initial-value problem is

\[
\theta(0)=\theta_0,\qquad \omega(0)=\theta'(0)=\omega_0.
\]

\subsection{Small-Angle Solution}\label{small-angle-solution}

For small angles measured in radians,

\[
\sin\theta\approx\theta.
\]

The nonlinear equation becomes

\[
\theta''+\frac gL\theta=0.
\]

Let \(\omega_* =\sqrt{g/L}\). The general solution is

\[
\theta(t)=A\cos(\omega_*t)+B\sin(\omega_*t).
\]

At \(t=0\),

\[
\theta(0)=A=\theta_0.
\]

Differentiating,

\[
\theta'(t)=-A\omega_*\sin(\omega_*t)+B\omega_*\cos(\omega_*t),
\]

so

\[
\theta'(0)=B\omega_*=\omega_0,
\qquad
B=\omega_0\sqrt{\frac Lg}.
\]

Therefore,

\[
\boxed{\theta(t)=\theta_0\cos\left(\sqrt{\frac gL}t\right)
+\omega_0\sqrt{\frac Lg}\sin\left(\sqrt{\frac gL}t\right)}.
\]

Writing this as \(C\cos(\omega_*t-\phi)\) gives

\[
C=\sqrt{\theta_0^2+\omega_0^2\frac Lg},
\qquad
\tan\phi=\frac{\omega_0}{\theta_0}\sqrt{\frac Lg}.
\]

The period satisfies \(\omega_*T=2\pi\), hence

\[
\boxed{T_0=2\pi\sqrt{\frac Lg}}.
\]

This small-angle period is independent of amplitude and initial angular velocity.

\subsection{Large-Angle Period}\label{large-angle-period}

For large angles, retain \(\sin\theta\). Conservation of energy gives, for release from rest at amplitude \(\theta_0\),

\[
\frac12mL^2\dot\theta^2+mgL(1-\cos\theta)
=mgL(1-\cos\theta_0).
\]

Using \(1-\cos\theta=2\sin^2(\theta/2)\),

\[
\dot\theta^2=\frac{4g}{L}
\left[\sin^2\left(\frac{\theta_0}{2}\right)-
\sin^2\left(\frac\theta2\right)\right].
\]

Set \(k=\sin(\theta_0/2)\) and use the substitution

\[
\sin\frac\theta2=k\sin\phi.
\]

After simplifying the quarter-period integral,

\[
\frac T4=\sqrt{\frac Lg}\int_0^{\pi/2}
\frac{d\phi}{\sqrt{1-k^2\sin^2\phi}}.
\]

Thus

\[
\boxed{T=4\sqrt{\frac Lg}K(k)},
\]

where \(K(k)\) is the complete elliptic integral of the first kind. In software such as SciPy, \passthrough{\lstinline!ellipk!} commonly takes the parameter \(m=k^2\):

\begin{lstlisting}[language=Python]
from scipy.special import ellipk
import numpy as np

def exact_period(L, g, theta0):
    m = np.sin(theta0 / 2.0) ** 2
    return 4.0 * np.sqrt(L / g) * ellipk(m)
\end{lstlisting}

\subsection{First-Order System}\label{first-order-system}

Define angular velocity \(\omega=\theta'\). Then

\[
\frac d{dt}\begin{bmatrix}\theta\\\omega\end{bmatrix}
=\begin{bmatrix}\omega\\-(g/L)\sin\theta\end{bmatrix}
=f(t,y),
\qquad
y(0)=\begin{bmatrix}\theta_0\\\omega_0\end{bmatrix}.
\]

RK4 is applied to this vector equation component by component.

\subsection{RK4 Pendulum Simulation}\label{rk4-pendulum-simulation}

\begin{lstlisting}[language=Python]
import numpy as np

def pendulum_derivatives(t, y, g, L):
    theta, omega = y
    return np.array([omega, -(g / L) * np.sin(theta)])

def rk4_step(f, t, y, h, g, L):
    k1 = f(t, y, g, L)
    k2 = f(t + h / 2, y + h * k1 / 2, g, L)
    k3 = f(t + h / 2, y + h * k2 / 2, g, L)
    k4 = f(t + h, y + h * k3, g, L)
    return y + h * (k1 + 2*k2 + 2*k3 + k4) / 6

def simulate(theta0=np.pi / 4, L=1.0, g=9.81,
             t_end=10.0, h=0.05):
    times = np.arange(0.0, t_end, h)
    states = np.zeros((len(times), 2))
    states[0] = [theta0, 0.0]
    for i in range(1, len(times)):
        states[i] = rk4_step(
            pendulum_derivatives, times[i - 1], states[i - 1], h, g, L
        )
    return times, states
\end{lstlisting}

The trajectory can be plotted as \(\theta(t)\), and the phase portrait is obtained by plotting angular velocity against angle.

\subsection{Exact Trajectory and Convergence}\label{exact-trajectory-and-convergence}

For release from rest, the nonlinear trajectory can be expressed with Jacobi elliptic functions. With

\[
k=\sin(\theta_0/2),\qquad m=k^2,\qquad u=\sqrt{g/L}\,t,
\]

one representation is

\[
\theta(t)=2\arcsin\left(k\frac{\operatorname{cn}(u,m)}{\operatorname{dn}(u,m)}\right).
\]

The numerical error at the time nodes is

\[
E_i=|\theta_{\mathrm{RK4}}(t_i)-\theta_{\mathrm{exact}}(t_i)|.
\]

Repeat the experiment with several step sizes, for example

\[
h\in\{0.1,0.05,0.025,0.0125,0.00625\}.
\]

If RK4 is in its asymptotic regime, halving \(h\) reduces the global error by approximately \(2^4=16\). A log-log graph of error versus \(h\) has slope approximately four:

\[
|E(h)|\approx Ch^4.
\]

\subsection{Period Extraction by Zero Crossings}\label{period-extraction-by-zero-crossings}

A numerical period can be obtained from successive turning points, where \(\omega(t)=0\). If \(\omega_n\) and \(\omega_{n+1}\) have opposite signs, linearly interpolate the zero:

\[
t_{\mathrm{root}}=t_n-\omega_n\frac{t_{n+1}-t_n}{\omega_{n+1}-\omega_n}.
\]

The time between every second turning point is one full oscillation period. Comparing this value with the elliptic-integral period measures the combined integration and event-location error.

\subsection{Symplectic Euler}\label{symplectic-euler}

The pendulum is conservative, so long-time energy behavior matters. Explicit RK4 is accurate over a fixed interval but may slowly distort energy over very long integrations. A semi-implicit Euler method updates velocity first and then angle:

\[
\begin{aligned}
\omega_{n+1}&=\omega_n-h\frac gL\sin\theta_n,\\
\theta_{n+1}&=\theta_n+h\omega_{n+1}.
\end{aligned}
\]

It is only first-order accurate, but its geometric structure often prevents the systematic energy drift seen in ordinary explicit Euler.

\section{Higher-Order Runge--Kutta Methods}\label{higher-order-rungekutta-methods}

An explicit Runge--Kutta method obtains several stage slopes and combines them with weights. Increasing the order requires satisfying increasingly many nonlinear order conditions. These conditions can be organized using rooted trees and B-series.

For explicit methods of order \(p\ge5\), the Butcher barrier implies that more than \(p\) stages are required. For example, fifth order needs at least six stages and sixth order needs at least seven stages. The additional evaluations reduce the efficiency advantage of a fixed high-order method.

Embedded methods such as RKF45 and Dormand--Prince use the same stages with two weight vectors \(b_i\) and \(b_i^*\):

\[
y_{n+1}=y_n+h\sum_i b_iK_i,
\qquad
\widehat y_{n+1}=y_n+h\sum_i b_i^*K_i.
\]

The difference

\[
\|y_{n+1}-\widehat y_{n+1}\|
\]

estimates the local error. The solver can then increase \(h\) when the estimate is small and decrease \(h\) when it is large. This produces adaptive step-size control without evaluating the differential equation twice with two unrelated methods.

\section{Practical Checklist}\label{practical-checklist}

\begin{itemize}
\tightlist
\item
  Check that interpolation nodes are distinct.
\item
  Use the interpolation error formula to identify smoothness and node-placement effects.
\item
  Distinguish local truncation error from accumulated global error.
\item
  Verify an observed convergence order with a refinement table or log-log plot.
\item
  For ODEs, check both accuracy and stability; a small local error does not guarantee a useful long-time solution.
\item
  For conservative dynamics, monitor energy and phase-space behavior.
\item
  Keep angle units in radians when evaluating trigonometric functions.
\item
  When using \passthrough{\lstinline!scipy.special.ellipk!}, pass the parameter \(m=k^2\), not the modulus \(k\) itself.
\end{itemize}
\end{document}
